\documentclass[12pt,a4paper]{article}
\usepackage[utf8]{inputenc}
\usepackage[spanish]{babel}
\usepackage{amsmath}
\usepackage{amssymb}
\usepackage{listings}
\usepackage{xcolor}
\usepackage{geometry}
\usepackage{hyperref}

\geometry{top=2.5cm, bottom=2.5cm, left=2.5cm, right=2.5cm}

% Configuración de colores profesionales
\definecolor{darkblue}{rgb}{0.0, 0.2, 0.4}
\definecolor{codegreen}{rgb}{0,0.6,0}
\definecolor{codegray}{rgb}{0.5,0.5,0.5}
\definecolor{codepurple}{rgb}{0.58,0,0.82}
\definecolor{backcolour}{rgb}{0.97,0.97,0.96}

\hypersetup{
    colorlinks=true,
    linkcolor=darkblue,
    urlcolor=cyan,
}

% Configuración del estilo de código para Python
\lstdefinestyle{mystyle}{
    backgroundcolor=\color{backcolour},   
    commentstyle=\color{codegreen}\itshape,
    keywordstyle=\color{darkblue}\bfseries,
    numberstyle=\tiny\color{codegray},
    stringstyle=\color{codepurple},
    basicstyle=\ttfamily\small,
    breakatwhitespace=false,         
    breaklines=true,                 
    captionpos=b,                    
    keepspaces=true,                 
    numbers=left,                    
    numbersep=5pt,                  
    showspaces=false,                
    showstringspaces=false,
    showtabs=false,                  
    tabsize=4,
    extendedchars=true,
    literate={á}{{\'a}}1 {é}{{\'e}}1 {í}{{\'i}}1 {ó}{{\'o}}1 {ú}{{\'u}}1 {ñ}{{\~n}}1 {Á}{{\'A}}1 {É}{{\'E}}1 {Í}{{\'I}}1 {Ó}{{\'O}}1 {Ú}{{\'U}}1 {Ñ}{{\~N}}1 {¿}{{\textquestiondown}}1 {¡}{{\textexclamdown}}1
}

\lstset{style=mystyle}

\title{\textbf{\color{darkblue}Redes Neuronales de Grafos (GNN): Convoluciones Topológicas}}
\author{Documentación Técnica de Inteligencia Artificial Industrial}
\date{\today}

\begin{document}

\maketitle

\noindent Las Redes Neuronales de Grafos (GNN) representan el punto exacto donde la topología matemática clásica se fusiona con los modelos predictivos multidimensionales.

Cuando aplicamos \textit{Deep Learning} para el modelado de diagnósticos en infraestructuras físicas (como predecir el desgaste de rodamientos en una batería de bombas), los modelos tradicionales (como las redes neuronales densas o CNNs) fallan porque asumen que los datos son independientes o están en una cuadrícula perfecta.

Las GNNs, y específicamente las Redes Convolucionales de Grafos (GCN), entienden que si un motor operando a 460V comienza a vibrar anómalamente, esa vibración se propagará mecánicamente a través de las tuberías hacia los equipos adyacentes. La topología dicta la influencia.

\section{Desglosando la Ecuación Matemática (Paso de Mensajes)}

La ecuación que tienes en tu documento captura elegantemente cómo la información fluye por la red en una sola capa de la red neuronal:

\begin{equation}
    H^{(l+1)} = \sigma \left( \tilde{D}^{-\frac{1}{2}} \tilde{A} \tilde{D}^{-\frac{1}{2}} H^{(l)} W^{(l)} \right)
\end{equation}

Vamos a diseccionarla pieza por pieza para entender su lógica geométrica:

\begin{itemize}
    \item \textbf{$H^{(l)}$ (Matriz de Características):} Es el estado actual de tus nodos. Si tienes $5$ motores y estás midiendo $2$ variables (temperatura y vibración), esta es una matriz de $5 \times 2$.
    \item \textbf{$W^{(l)}$ (Matriz de Pesos):} Son los parámetros entrenables. Es la parte que la inteligencia artificial ajusta durante el entrenamiento para aprender qué características son realmente importantes para predecir un fallo.
    \item \textbf{$\tilde{A}$ (Matriz de Adyacencia con Autoenlaces):} Normalmente, la matriz de adyacencia $A$ solo te dice quién está conectado con quién. Pero al actualizar el estado de un equipo, no solo te importa lo que le pasa a sus vecinos, ¡también te importa su propio estado actual! Por eso se le suma la matriz identidad ($\tilde{A} = A + I$).
    \item \textbf{$\tilde{D}^{-\frac{1}{2}} \dots \tilde{D}^{-\frac{1}{2}}$ (Normalización Simétrica):} Este es el secreto matemático para que la red no colapse. Si un colector principal está conectado a 10 bombas, al sumar la información de todas ellas, sus valores numéricos explotarían en comparación con una bomba aislada. Dividir por la raíz cuadrada de los grados ($\tilde{D}$) de ambos nodos conectados asegura que el flujo de información se normalice de manera justa, amortiguando la influencia de los nodos masivamente conectados.
    \item \textbf{$\sigma$ (Activación):} Una función no lineal (como ReLU) que permite a la red aprender patrones complejos, no solo relaciones lineales simples.
\end{itemize}

\section{Implementación en Python: Construyendo la Ecuación 7 desde Cero}

Aunque librerías como \textit{PyTorch Geometric} hacen esto con una sola línea de código en producción, la mejor forma de entender la arquitectura profunda es programar la matriz exacta usando tensores nativos de PyTorch.

Vamos a simular un paso de mensaje convolucional en una pequeña red de tres equipos interconectados.

\begin{lstlisting}[language=Python, caption={Paso de mensajes en una GCN usando PyTorch}]
import torch
import torch.nn as nn
import numpy as np

# --- 1. Definicion de la Topologia ---
# Grafo de 3 nodos (0: Tablero, 1: Motor A, 2: Motor B)
# Conexiones: Tablero <-> Motor A, Motor A <-> Motor B
A = torch.tensor([[0., 1., 0.],
                  [1., 0., 1.],
                  [0., 1., 0.]])

# Anadimos auto-enlaces (Matriz Identidad) para que cada nodo recuerde su propio estado
I = torch.eye(3)
A_tilde = A + I

# --- 2. Normalizacion Simetrica ---
# Calculamos el grado de cada nodo (cuantas conexiones tiene, incluyendo el autoenlace)
grados = A_tilde.sum(dim=1)

# Matriz diagonal D_tilde^(-1/2)
D_tilde_inv_sqrt = torch.diag(torch.pow(grados, -0.5))

# Calculamos el operador de paso de mensajes: D^(-1/2) * A_tilde * D^(-1/2)
operador_convolucional = torch.matmul(torch.matmul(D_tilde_inv_sqrt, A_tilde), D_tilde_inv_sqrt)

print("Operador Convolucional Normalizado:")
print(np.round(operador_convolucional.numpy(), 3))

# --- 3. Datos de los Nodos (Matriz H) ---
# Supongamos 2 caracteristicas por equipo: [Temperatura, Nivel de Vibracion]
H_l = torch.tensor([[45.0, 0.1],  # Tablero (frio, sin vibracion)
                    [80.0, 2.5],  # Motor A (caliente, vibrando mucho)
                    [50.0, 0.5]]) # Motor B (normal)

# --- 4. Parametros Entrenables (Matriz W) ---
# Proyectamos de 2 caracteristicas a 3 caracteristicas latentes (ocultas)
W_l = nn.Parameter(torch.rand(2, 3))

# --- 5. APLICACION DE LA ECUACION (7) ---
# H^(l+1) = sigma( D^(-1/2) A_tilde D^(-1/2) H^(l) W^(l) )

# a. Paso de mensajes (Los nodos mezclan su estado con sus vecinos)
mezcla_vecinal = torch.matmul(operador_convolucional, H_l)

# b. Transformacion lineal (Aprendizaje)
transformacion = torch.matmul(mezcla_vecinal, W_l)

# c. Activacion no lineal (sigma)
H_l_mas_1 = torch.relu(transformacion)

print("\nMatriz de Caracteristicas de Entrada H^(l):")
print(H_l)
print("\nNuevo Estado Oculto de los Nodos H^(l+1):")
print(H_l_mas_1.detach().numpy())
\end{lstlisting}

\section{¿Qué está sucediendo aquí físicamente?}

Si observas la \texttt{mezcla\_vecinal} antes de la transformación lineal, verías que el Motor B (que originalmente estaba a $50^\circ$C y $0.5$ de vibración) absorbió matemáticamente parte de las características del Motor A. La red neuronal acaba de modelar la transferencia térmica y armónica a través del acople estructural.

Al apilar múltiples capas (repitiendo esta ecuación), el nodo $0$ recibe información del nodo $2$ a través del nodo $1$. Es decir, con $k$ capas convolucionales, cada nodo adquiere ``conciencia'' geométrica de todos los componentes que están a $k$ saltos de distancia.

Para consolidar esta abstracción matemática, he diseñado un simulador donde puedes observar cómo los valores numéricos se promedian y propagan por la red topológica en cada iteración del paso de mensajes.

\end{document}
